\documentclass[11pt]{article}
\usepackage[a4paper,margin=1.6cm]{geometry}
\usepackage{fontspec}
\usepackage{array}
\setmainfont{LiberationSerif}[
  Path = /sessions/lucid-sleepy-darwin/mnt/Vietnam/MedicalRecord/handwriting/fonts/, Extension = .ttf,
  UprightFont = *-Regular, BoldFont = *-Bold,
  ItalicFont = *-Italic, BoldItalicFont = *-BoldItalic,
]
\newfontfamily\hw{PlaywriteVN}[Path=/sessions/lucid-sleepy-darwin/mnt/Vietnam/MedicalRecord/handwriting/fonts/, Extension=.ttf, UprightFont=*]
\renewcommand{\arraystretch}{1.6}
\pagestyle{empty}
\begin{document}
\footnotesize
\noindent
\begin{minipage}[t]{0.55\linewidth}
Bộ Y tế/Sở Y tế: {\hw Sở Y tế Đà Nẵng} \\
Bệnh viện: Bệnh viện Phục hồi chức năng \\
Khoa/Trung tâm: Phục hồi chức năng \\
Buồng {\hw 4} \quad Giường {\hw 2}
\end{minipage}
\hfill
\begin{minipage}[t]{0.4\linewidth}
\raggedleft
(Mẫu phiếu số 3) \\
MS: 27/BV1-P3 \\
Số vào viện 000123
\end{minipage}

\begin{center}
\textbf{\large PHIẾU THỰC HIỆN KỸ THUẬT} \\
\textbf{\large PHỤC HỒI CHỨC NĂNG}
\end{center}

\noindent
- Họ tên người bệnh: {\hw Lê Văn Hòa}
  \quad Năm sinh/Tuổi: {\hw 61}
  \quad Nam/Nữ: {\hw Nam} \\
- Khoa: Phục hồi chức năng \quad Số buồng: {\hw 4}
  \quad Số giường: {\hw 2} \\
- Chẩn đoán: {\hw Chấn thương tủy sống ngực D10 liệt hai chi dưới}

\vspace{0.3cm}
\begin{tabular}{|p{1.7cm}|p{3.9cm}|p{3.7cm}|p{1.0cm}|p{1.9cm}|p{2.1cm}|p{1.5cm}|}
\hline
\textbf{Ngày giờ} &
\textbf{Mô tả diễn biến} &
\textbf{Tên dịch vụ kỹ thuật} &
\textbf{TG (phút)} &
\textbf{Người thực hiện} &
\textbf{BS chỉ định} &
\textbf{Xác nhận} \\
\hline
{\hw 14h00 13/06} & {\hw Biên độ gấp khuỷu cải thiện, chưa đau khi vận động} & {\hw Tập dưỡng sinh} & {\hw 15} & {\hw Hoàng Văn Tùng} & {\hw BS. Trần Thị Bích} & {\hw X} \\ \hline
{\hw 14h45 13/06} & {\hw Phù nề giảm, tiếp tục dẫn lưu bạch huyết} & {\hw Kỹ thuật kiểm soát tư thế (ngồi, bò, đứng, đi)} & {\hw 15} & {\hw Hoàng Văn Tùng} & {\hw BS. Lê Văn Cường} & {\hw Đã ký} \\ \hline
{\hw 14h30 13/06} & {\hw Cơ lực chi dưới cải thiện, còn hạn chế duỗi gối} & {\hw Kỹ thuật tập sức bền có gắn theo dõi tim mạch} & {\hw 15} & {\hw Nguyễn Thị Hoa} & {\hw BS. Lê Văn Cường} & {\hw Đã ký} \\ \hline
{\hw 14h15 13/06} & {\hw Bệnh nhân hợp tác tốt, mức độ đau giảm rõ} & {\hw Kỹ thuật điều trị bằng máy ép khí ngắt quãng} & {\hw 45} & {\hw Hoàng Văn Tùng} & {\hw BS. Nguyễn Văn An} & {\hw X} \\ \hline
{\hw 14h30 14/06} & {\hw Bệnh nhân mệt sau 15 phút tập, giảm cường độ buổi tới} & {\hw Kỹ thuật thư giãn} & {\hw 20} & {\hw Hoàng Văn Tùng} & {\hw BS. Nguyễn Văn An} & {\hw Đã ký} \\ \hline
{\hw 14h00 14/06} & {\hw Khả năng tự xoay trở trên giường tốt hơn} & {\hw Điều trị bằng laser công suất thấp chiếu ngoài} & {\hw 30} & {\hw Lê Thị Mai} & {\hw BS. Phạm Thị Dung} & {\hw X} \\ \hline
\end{tabular}

\end{document}
